% Canonical statement: sections/05_learners.tex:6-21
\begin{theorem}[Pointwise cap mixture and table-polynomial implementation]\label{thm:caps}
Fix $d\ge1$. Let $F$ have a strict rank-$d$ representation and no negative $2\times2$, and let $A$ be any monotone flip of its negative entries. Perturb the row vectors once into linear general position without changing $F$. At any nonempty survivor set, use the uniform row prior $\nu$. There is a cap sampler, depending on the surviving rows but not on the unknown marginal, such that
\begin{equation}\label{eq:capcoverage}
 \forall x\in U,\qquad \E[\1_S(x)\nu(T)]\ge\frac{c_d}{8d},
 \quad c_d=\Pr_{u\sim\mathrm{Unif}(S^{d-1})}
 [u_1\ge\sqrt{1-1/(2d)}].
\end{equation}
Here $c_1=1/2$. Thus Theorem~\ref{thm:main} applies with $\rho=8d/c_d$ in the ideal real-arithmetic model.

For each fixed $d$, given the explicit table of $A$ and rational template vectors of bit length at most $b$, there are randomized and deterministic implementations, polynomial in $|X|,|H|,b,1/\epsilon$, for every $0<\epsilon<1/4$, with
\[
 m=O_d(\log|H|+\log(1/\epsilon)),\qquad
 \tau=\Omega_d(\epsilon/\log(1/\epsilon)).
\]
Queries are intersections of stored cap tests, or such tests split by one row label or one point's row labels. The running time is polynomial in the supplied table size.
\end{theorem}

% Proof binding: appendices/cap_proper_proofs.tex:4-57
\begin{proof}[Proof of \cref{thm:caps}]
We first make the once-only perturbation precise. Number the rows $j=1,\ldots,K$. Replace a row vector $v_j$ by
\[
 v_j(t)=v_j+t(1,j,\ldots,j^{d-1}).
\]
Choose a positive rational upper bound $t_0$ so that every dot product retains its strict sign for $0<t<t_0$. The finitely many original rational dot products and the perturbation dot products give such a bound by direct comparison. For each $d$-subset of rows, its determinant is a polynomial in $t$ of degree $d$ whose leading coefficient is a nonzero Vandermonde determinant. At most $d\binom Kd$ values are forbidden. Testing that many plus one rational grid values in $(0,t_0)$ finds a valid value, with polynomial bit length and polynomial work for fixed $d$. When $K<d$, no perturbation is needed: use the evaluation-coordinate construction below for every survivor set. For $K\ge d$, the same perturbation works for every later survivor set.

Suppose there are $k>d$ surviving rows. Any proper $j$-dimensional subspace contains at most $j$ of them, strictly less than $kj/d$. The radial-isotropy proof of Theorem~\ref{thm:forster} therefore applies; its convex formulation is also proved in Appendix~\ref{app:convex}. Transform and normalize the rows to vectors $v_h\in S^{d-1}$ with
\[
 \E_\nu v_hv_h^\top=I_d/d.
\]
Transform points contragrediently and normalize them to $z_x$. All template signs remain strict. If $k\le d$, retain \emph{all} rows and replace the factorization by evaluation coordinates: let $\widetilde z_x=(\langle v_h,u_x\rangle)_{h\in V}$ and use row vectors $e_h\in\R^k$. Every coordinate of $\widetilde z_x$ is nonzero by strictness, so its normalization is defined and preserves every sign. For the uniform surviving-row prior, $k^{-1}\sum_{h\in V}e_he_h^\top=I_k/k$. This is a new factorization, not an invertible map of the original ambient space. No target is discarded.

Write the effective dimension as $j\le d$, put $\alpha^2=1/(2j)$ and $\beta=\sqrt{1-\alpha^2}$, and sample $u$ uniformly on $S^{j-1}$. Define
\[
 C_u=\{x:\langle u,z_x\rangle\ge\beta\},\qquad
 T_u^\pm=\{h:\pm\langle u,v_h\rangle>\alpha\}.
\]
For $Z=\langle u,v_h\rangle^2\in[0,1]$, isotropy gives $\E_\nu Z=1/j$. Since
\[
 \E Z\le\alpha^2+\Pr[Z>\alpha^2],
\]
we have $\nu(T_u^+)+\nu(T_u^-)\ge1/(2j)$. Choose the larger side, of mass at least $1/(4j)$.

The geometry is strict. If $s=\langle u,v\rangle>\alpha$ and $t=\langle u,z\rangle\ge\beta$, then
\[
 \langle z,v\rangle\ge ts-\sqrt{1-t^2}\sqrt{1-s^2}>0,
\]
because $s,t>0$ and $s^2+t^2>1$. Replacing $v$ by $-v$ proves the negative-side assertion. Thus $C_u\times T_u^+$ is positive in $F$, and $C_u\times T_u^-$ is negative in $F$.

A positive rectangle survives in $A$. A negative rectangle has a singleton side. If its row side is $\{h\}$, partition $C_u\cap U$ according to $h(x)$ in $A$ and choose either color with probability $1/2$. If its point side is $\{x\}$, partition the row side according to $h(x)$ and choose either color with probability $1/2$. Empty outcomes are allowed. The singleton-point case is essential and is not replaced by a singleton-row assumption. Conditional on $x\in C_u$, its expected weighted coverage after this split is at least half the chosen side's mass, hence at least $1/(8j)$. Rotational invariance gives probability $c_j$ of cap membership. Therefore the pointwise bound is $c_j/(8j)\ge c_d/(8d)$. To verify the last inequality, couple the first coordinate of uniform spheres by normalizing an increasing list of independent standard Gaussians. Its absolute value decreases as dimension grows, while $\sqrt{1-1/(2j)}$ increases. Symmetry supplies the factor $1/2$, including the case $j=1$.

For finite-bit computation, Appendix~\ref{app:convex} gives a deterministic rational transform with normalized covariance at least $3I_j/(4j)$. The preceding second-moment argument then gives total side mass at least $1/(4j)$, selected-side mass at least $1/(8j)$, and split coverage at least $1/(16j)$ whenever a point belongs to the cap.

Use the finite direction set
\[
 W_j=\left\{\frac{g}{\|g\|}:0\ne g\in(16j)^{-1}\mathbb Z^j\cap[-1,1]^j\right\},
 \qquad J_j=|W_j|\le(32j+1)^j.
\]
Repetitions of a direction cause no difficulty. Rounding a unit vector coordinatewise produces $g$ with $\|g-z\|\le1/(32\sqrt j)$; normalization changes it by at most twice this amount. Thus $\langle z,g/\|g\|\rangle\ge1-1/(512j)>\beta$, where the last inequality follows from $\beta\le1-1/(4j)$. Every point is in at least one cap. Uniformly sampling $W_j$ and the split color yields coverage at least $1/(16jJ_j)$. The common rational lower bound
\begin{equation}\label{eq:netr}
 r_d=\frac1{16d(32d+1)^d}
\end{equation}
works in all effective dimensions. All membership comparisons are rational: for example $\langle z,g/\|g\|\rangle\ge\beta$ means that the unnormalized dot product is positive and
\[
 \langle z_{\rm raw},g\rangle^2\ge(1-1/(2j))\|z_{\rm raw}\|^2\|g\|^2.
\] The row-side tests use the analogous strict squared comparison. Consequently no exact sampling of a continuous direction is required. Theorem~\ref{thm:main} with $r_d$ proves the randomized finite-bit bound.

For determinism, enumerate the at most $2J_j$ outcomes and query each $D(S)$, using the same tolerance and stopping check as Theorem~\ref{thm:main} with $r_d$. Choose the outcome maximizing $b\widehat s$, where $b=\nu(T)$. Coverage implies that some outcome has $bs\ge r_du$. The chosen one has
\[
 bs\ge r_du-2\tau,\qquad ab\ge r_d-2\tau/u\ge3r_d/4.
\]
Thus both $a$ and $b$ are at least $3r_d/4$. Query its mismatch mass and peel or eliminate exactly as before. The target survives, and all error bounds remain valid. Each stage has progress at least $3r_d/4$, whereas the deterministic total-progress bound is $B$. There cannot be more than $\lceil4B/(3r_d)\rceil+1$ nonstopping stages. There are $O_d(1)$ queries per stage. This proves the deterministic claim, without fixing a successful private seed after seeing $D$.
\end{proof}
